\newcommand{\AngryAny}{11.95}
\newcommand{\AngryRR}{1.809}
\newcommand{\TailNetPercent}{76.34\%}
\newcommand{\TailNetValue}{-90.04}
\newcommand{\BodyNetValue}{-27.91}
\newcommand{\AngryAnyCI}{[7.89, 16.01]}
\newcommand{\AngryRRCI}{[1.316, 2.488]}
\newcommand{\SharesAny}{2.11}
\newcommand{\SharesAnyCI}{[0.90, 3.31]}
\newcommand{\SharesRR}{1.016}
\newcommand{\SharesRRCI}{[0.826, 1.250]}
\newcommand{\SadAny}{7.57}
\newcommand{\SadAnyCI}{[3.78, 11.35]}
\newcommand{\SadRR}{0.723}
\newcommand{\SadRRCI}{[0.464, 1.126]}
\newcommand{\CellNet}{-100.39}
\newcommand{\CellNetCI}{[-142.75, -58.03]}
\newcommand{\CellTheta}{1.282}
\newcommand{\CalibrationResults}{In the full-size design-anchored source-plus-block simulation, family noncoverage is 30.0\% for PPML and 1.0\% for occurrence projections, with 100 replicates and wide Monte Carlo intervals. No simultaneous positive/negative count declaration occurs in those 100 realizations. Table~\ref{tab:calibration} reports invalid cases and the smaller-design scenarios. The reference parameters marginalize over re-drawn source/block slopes; these are not error bounds for a finite-source conditional parameter or estimates of how often the observed four-sign profile is false. The failure of nominal count-interval coverage limits their inferential interpretation even when fitted mean profiles remain stable.}
